% Fragmento público generado desde propietarios íntegros.
% Residencia: 52.5.2.
% Fuente: ../incoming/radion_monografia_20260827/manuscrito/sections/06_interior_helice.tex:109-201
\subsubsection{Relèvement hélicoïdal du retour nonadique}

\paragraph{Décomposition nonadique et quotient de mémoire}

L'opération élémentaire
\begin{equation}
n=9q_9(n)+\rho_9(n),
\qquad
H_9(n)=(\rho_9(n),q_9(n))
\label{eq:H9}
\end{equation}
satisfait
\begin{equation}
H_9(n+9)=H_9(n)+(0,1).
\label{eq:H9-helix}
\end{equation}
La première coordonnée se ferme ; la seconde s'élève. Cette identité est la
version arithmétique minimale d'une vis.

Pour une fibre spectrale,
\[
n\longmapsto(\tau_\gamma^n,n)
\]
est une hélice sur \(S^1\). Ajouter la contraction \((5/9)^n\) produit une
spirale intérieure. Le radion participe aux deux lectures : une unité de phase
à la frontière et une unité orientée d'action dans le relèvement.

\paragraph{Rang de mémoire comme lecteur dimensionnel}

Supposons qu'une projection visible conserve la valeur de phase, mais perde
\(d\) cocycles de mémoire linéairement indépendants. Tout relèvement fidèle
nécessite au moins \(d\) coordonnées supplémentaires pour les reconstruire. Cela
motive une lecture précise de la mécanique dimensionnelle :
\begin{equation}
\begin{aligned}
\text{dimension émergente}
&=\text{rang minimal de mémoire indépendante}\\
&\phantom{=}\text{non publiable sur le feuillet visible}.
\end{aligned}
\label{eq:dimension-memory}
\end{equation}
On n'affirme pas que toute dimension physique soit un compteur. On affirme que, dans le
domaine TPK, chaque mémoire indépendante qui doit survivre impose une
coordonnée du relèvement.

\begin{narrativebox}
La dimension apparaît lorsque la projection ne peut plus conserver toute
l'histoire. Le cercle suffit à dire où se trouve la phase ; il ne
suffit pas à dire combien de fois elle est revenue, quel feuillet elle a transporté, quel intervalle
s'est contracté ou quelle retenue a survécu.
\end{narrativebox}

\paragraph{Assemblage de cylindres, de spirales et de trajectoires hélicoïdales}

Les cylindres coinductifs ne sont pas des tubes physiques ajoutés à une spirale. Ce sont
des intervalles emboîtés qui conservent les préfixes et resserrent la valeur. La spirale
décrit la contraction transversale ; l'hélice décrit la mémoire longitudinale ;
le cylindre décrit l'ensemble des valeurs compatibles à une profondeur
finie. À la limite, l'intersection des cylindres sélectionne une section
unique.

L'image spatiale réunit, sans confondre :
\[
\begin{array}{ccl}
\text{angle} &\leftrightarrow& \text{phase visible},\\
\text{hauteur} &\leftrightarrow& \text{retour/mémoire},\\
\text{rayon de cylindre} &\leftrightarrow& \text{incertitude de raffinement},\\
\text{feuillet} &\leftrightarrow& \text{orientation bidirectionnelle}.
\end{array}
\]
Le taux \(3^{-54}\) contracte le cylindre à chaque retour ; il ne s'ajoute pas à l'angle et
ne remplace pas la retenue réelle \(\epsR\).

\subsubsection{Clôtures de phase d'ordres \(2\pi\) et \(4\pi\)}

Dans la projection circulaire, \(2\pi\) restaure la phase. Dans un relèvement orienté
à deux feuillets, un tour peut se fermer sur le feuillet conjugué et nécessiter un
second tour pour restaurer l'orientation :
\[
U_{2\pi}=-I,\qquad U_{4\pi}=I.
\]
La structure coïncide avec la périodicité spinorielle lorsque l'on applique le
foncteur à une représentation \(SU(2)\) ; on n'identifie pas le relèvement TPK à
\(SU(2)\) sans cette application.

La lecture surfacique est compatible :
\[
\mathcal A(2\pi)=\hfull,\qquad
\mathcal A(4\pi)=2\hfull.
\]
Dans un ruban de Möbius, un circuit visible inverse le feuillet et deux circuits
restaurent l'orientation. Le \(2\pi/4\pi\) n'est pas une duplication accidentelle ;
c'est la différence entre retour de phase et retour de l'état enrichi.
